% Shoggoth 21 Aug 26 

\subsection{On “Lay Activity”}

Editorial article in “Foldebladet” for April 3, 1907.

The large Joint Committee which is negotiating concerning the union of the church bodies also dealt with “lay activity” at its meeting last autumn. The result of the work is set forth in seven theses which, it is said, were unanimously adopted by the representatives of the different bodies. To each of the seven theses there has been appended one or more passages of Scripture, or else statements from the Lutheran Confessional Writings. The connection between the Scripture passages and the theses is, however, so unclear that interpretation is needed in order to make them fit together.

Nor do the theses speak of “lay activity” alone, not even chiefly. They speak of the congregation, the means of grace, the preaching office, the gifts of grace, the universal priesthood, and finally perhaps of “lay activity.” But although all these things are spoken of, there is nevertheless no explanation either of what is meant by the congregation or by “lay activity.” Indeed, we say that perhaps “lay activity” is spoken of; for it is not mentioned by name in any of the theses. Strictly speaking, we infer that Thesis 5 speaks of “lay activity,” because there are certain expressions which somewhat recall this matter, and because such great care is taken to have it properly controlled.

Proposition 5 reads as follows:

\begin{quote}
God also wills that the special gifts of grace which he may have given to certain persons in the congregation, and which therefore are the property of the congregation, shall be put to use by the congregation; and thus in particular the gift of grace of being able to proclaim God’s Word before an assembly. According to God’s Word, the one who has received this gift of grace shall be called to this service in the congregation by its request or consent, and shall use the gift under the oversight of the congregation or of those whom the congregation has called to exercise oversight over the work in the congregation, as God’s Word and our Church’s Confession say.
\end{quote}

Nothing is said about what is meant by the congregation, either in this thesis or in any of the others; we read only in Thesis 1 that “God has given to the congregation, also to the individual local congregation, the means of grace,” etc. Presumably, therefore, by the congregation is meant simply the outward organization which in each place bears the name congregation, whether it deserves the name or not, whether it is dead or living, awake or asleep.

If we look at the thesis on this assumption, we find the following assertions, some with and some without the assurance that they are “according to God’s Word”:

\begin{enumerate}

\item When certain persons in the congregation have received a gift of grace, that gift is therefore “the property of the congregation.”


\item When someone has received “the gift to proclaim God’s Word before an assembly,” he shall be called to this service in the congregation by its request or consent.


\item When he has thus been called, he shall use the gift under the oversight of the congregation or of those whom the congregation has called to exercise oversight over the work in the congregation.

\end{enumerate}

The first of these assertions stands on its own merits; but the second is said to be “according to God’s Word”; the third, “as God’s Word and our Church’s Confession teach.”

We have understood the Joint Committee in this Thesis 5 to be thinking of “lay activity”; for the quotation marks which are used must surely indicate here that the expression is taken from ordinary usage, which has made the term so well known among us; and the “lay activity” commonly known among us may perhaps be described as a “gift of grace to proclaim God’s Word in an assembly,” without our thereby acknowledging that the explanation is exact or sufficient.

Now it is asserted in Thesis 5 that all gifts of grace which God has given to certain persons in the congregation are therefore the property of the congregation. From what follows, the meaning is that the congregation can and shall determine whether these gifts are to be used and how they are to be used, of course under responsibility to God, but with such power and authority that no one has the right to use the gift—at least not the gift of speaking in an assembly—without the congregation’s permission.

Here we must ask: Where is it written that the individual Christian’s gift of grace is the property of the congregation? “All things are yours,” says Paul, but surely that is meant of the Christians, not of an outward organization of people concerning whose Christianity no more can be said than that they are entered in a congregational register? Christians can say that all things are ours because they are Christ’s and Christ is God’s; but surely it is not thereby said that an outward organization of Christians and non-Christians may dispose of and command all things in heaven and on earth, and in particular may determine whether the Christian shall use his gift of grace or not? Can a congregational decision remove the responsibility from those who have received the gift, so that a man need not use it because the congregation does not call upon him to do so? God’s Word says that the individual Christian bears responsibility for his gift of grace or his talent. God’s Word says that only the one who has gained a return from his pound enters into his master’s joy, and that the one who has no return, because he has buried his pound so carefully in the ground, is cast into the outer darkness, where there is weeping and gnashing of teeth. See, that speech has an altogether different sound from the Committee’s Thesis 5, and it gives no excuse to the lazy flesh.

No, one says, the Committee has meant nothing of the kind except with respect to the gift of proclaiming God’s Word in an assembly. “Lay activity” in the Haugean sense, then, is the one gift which the congregation must acknowledge by request or consent, and the one who has received it must conduct himself accordingly. If he receives a request, he shall speak; if he receives no request, he shall be silent in an assembly.

If this is the meaning, does it make the matter any better? Or is it written anywhere that the one who has received this gift is free of his responsibility when the congregation does not call upon him?

Suppose now that the congregation sleeps, and God in his grace sends a crying voice to awaken it? Or suppose that the congregation is dead, even though it has the name that it lives; must the one who was sent even to awaken the spiritually dead then be silent until the dead call upon the living one to speak?

No, neither the first assertion, that the gifts of grace are the property of the congregation, nor the second assertion, that no one may bear witness in an assembly before he has been called upon by the congregation, holds good “according to God’s Word,” which, on the contrary, makes the duty of witness dependent not upon a congregational call but simply and directly upon the personal Christian calling. Whether the witness is given to an individual or “in an assembly” makes no difference according to Scripture. Let the one who has the gift to speak to the individual speak to the individual, and let the one who has the gift to speak in an assembly speak in an assembly.

What is meant, after all, by “an assembly”? Notice that it does not say the congregation’s assembly, or the congregation; it says merely “an assembly.” How many constitute an assembly, and where and when may people come together before it becomes an assembly? Anyone can see that such indefinite words can be used by a misguided congregation or pastor to suppress the proclamation of God’s Word and every testimony of individual Christians both within and outside the congregation. French law in Madagascar provides that a Christian evangelist or missionary may, outside the church building, speak God’s Word to two but not to three, for three are an assembly; at home in France the law says that a Christian preacher may speak God’s Word outside the church building to nineteen but not to twenty persons, for twenty persons are an assembly. Is it not possible that the enemies of living Christianity are of the same kind in America as in France? And are there not many, many congregations in which the majority is opposed to lay activity?

We ought not frame our rules in such a way that the enemies of the awakening are given so useful a handle to seize upon as this. It was precisely by such means that Hauge was made a criminal in Norway, because without the call of men he spoke God’s Word in assemblies. Are we to return to such rules, which here in our country may indeed be unable to bring a man to fines and imprisonment, but which can nevertheless close his mouth, so that he is silent where he ought to speak?

And the final assertion in Thesis 5, that lay activity shall be exercised under the oversight of the congregation or of those whom the congregation has chosen for that purpose, plainly moves in the same direction. For in this context such oversight must mean that the congregation, or pastor, or congregational council has the right once again to deny the one who has the gift the use of it, if they so choose. For it comes to the same thing: God may give gift and call, but men, because they have formed an outward congregational order, may forbid the use of the gift.

To all this the Committee can say that it has a Thesis 6 which declares that in such a case “the congregation suffers harm.” But that does not meet the issue; for the question is this: Can the one who has received the gift refrain from using it for any reason whatever without himself suffering harm? And it is a harm which, according to God’s Word, consists in this, that “his hands and feet are bound, and he is cast into the outer darkness, where there is weeping and gnashing of teeth.”

[Georg Sverdrup.]

